% Irodov Part V §5.5 Dispersion and Absorption, original problems 200--215.
% Source PDF pp.229--232 (printed pp.220--223), answers pp.343--344.
% Answers and extracted equations are cross-checked against actual page images
% where radicals or coefficients can be lost during PDF text extraction.

\dispersionproblem{200}{220}
自由电子处于一束单色光波中，
其强度 $I=150\,\mathrm{W/m^2}$，
角频率 $\omega=3.4\times10^{15}\,\mathrm{s^{-1}}$。
求 (a) 电子受迫振动的位移、速度振幅；
(b) 光波磁场力与电场力\emph{振幅}之比，
并证明该比值为 $v_m/(2c)$。
\exsolution
取 $E(t)=E_m\cos\omega t$，
$I=\tfrac12\varepsilon_0cE_m^2$，得
$E_m=\sqrt{2I/(\varepsilon_0c)}\simeq336\,\mathrm{V/m}$。
忽略磁场力，电子满足
$m_e\ddot x=-eE_m\cos\omega t$，稳态振动的
\[
 \boxed{a=\frac{eE_m}{m_e\omega^2}
 \simeq5.1\times10^{-18}\,\mathrm m},\qquad
 \boxed{v_m=a\omega
 \simeq1.73\times10^{-2}\,\mathrm{m/s}}.
\]
由于 $B=E/c$，速度与电场相差 $\pi/2$，
磁力的时间因子为 $\sin\omega t\cos\omega t$，
其最大值含因子 $1/2$：
\[
 \boxed{\frac{F_{m,\max}}{F_{e,\max}}
 =\frac{v_m}{2c}\simeq2.9\times10^{-11}}.
\]
\excheck $a\simeq5.1\times10^{-16}\,\mathrm{cm}$；
磁场作用确实可在求一阶运动时忽略。

\dispersionproblem{201}{221}
角频率 $\omega$ 的电磁波在稀薄等离子体中传播，
自由电子数密度 $n_0$，
忽略离子的直接响应。求
(a) 相对介电常数 $\varepsilon_r(\omega)$；
(b) 相速度随\emph{介质内波长} $\lambda$ 的关系。
\exsolution
自由电子的位移满足
$-m_e\omega^2x=-eE$，
极化强度 $P=-n_0ex=-n_0e^2E/(m_e\omega^2)$。
定义
\[
 \omega_p^2=\frac{n_0e^2}{m_e\varepsilon_0},
 \qquad
 \boxed{\varepsilon_r=1-\frac{\omega_p^2}{\omega^2}}.
\]
波方程给 $\omega^2=c^2k^2+\omega_p^2$。
用介质内波长 $k=2\pi/\lambda$，得
\[
 \boxed{v_{\rm ph}=\frac\omega k
 =c\sqrt{1+\frac{\omega_p^2\lambda^2}
 {4\pi^2c^2}}}.
\]
\excheck $\omega>\omega_p$ 才有传播波；
在无碰撞稀薄等离子体中，相速度可大于 $c$，
这不代表信号速度超过 $c$。

\dispersionproblem{202}{221}
电离层对频率 $\nu=100\,\mathrm{MHz}$
的无线电波折射率为 $n=0.90$。
求自由电子数密度。
\exsolution
由等离子体
$n^2=1-\omega_p^2/\omega^2$，
$\omega=2\pi\nu$，得
\[
 \boxed{n_0=
 \frac{4\pi^2\nu^2m_e\varepsilon_0}{e^2}(1-n^2)}.
\]
代入基本常数，
\[
 \boxed{n_0\simeq2.4\times10^{13}\,\mathrm{m^{-3}}
 =2.4\times10^7\,\mathrm{cm^{-3}}}.
\]
\excheck 密度以 $\mathrm{cm^{-3}}$ 给出时须换算 $10^6$ 倍。

\dispersionproblem{203}{221}
硬 X 射线的真空波长为 $\lambda=50\,\mathrm{pm}$。
将石墨电子近似看作自由电子，
求石墨的折射率偏离1的数量。
\exsolution
电子等离子体色散在 $\omega\gg\omega_p$ 时给
$n=\sqrt{1-\omega_p^2/\omega^2}
\simeq1-\omega_p^2/(2\omega^2)$，
$\omega=2\pi c/\lambda$。
碳原子含6个电子，数密度
$n_e=6\rho N_A/M_C$，因而
\[
 \boxed{n-1\simeq-
 \frac{n_e e^2\lambda^2}{8\pi^2\varepsilon_0m_ec^2}}.
\]
用石墨密度约
$\rho=2.26\,\mathrm{g/cm^3}$、
$M_C=12.01\,\mathrm{g/mol}$，得
\[
 \boxed{n-1\simeq-5.4\times10^{-7}}.
\]
\excheck X 射线下 $n<1$，差值是无量纲且很小；
所用石墨密度为材料近似值。

\dispersionproblem{204}{221}
电子在谐振回复力 $-\kappa x$ 与阻尼力
$-\gamma\dot x$ 下，受交变电场
$E=E_0\cos\omega t$ 驱动。
忽略磁力，求 (a) 稳态位移；
(b) 平均吸收功率、最大吸收频率及最大值。
\exsolution
设电子电荷 $-e$，
$m_e\ddot x+\gamma\dot x+\kappa x=-eE_0\cos\omega t$。
记 $\omega_0^2=\kappa/m_e$，
$\beta=\gamma/(2m_e)$，
稳态可写作 $x=a\cos(\omega t+\phi)$，其中
\[
 \boxed{a=\frac{eE_0/m_e}
 {\sqrt{(\omega_0^2-\omega^2)^2+4\beta^2\omega^2}}},
 \quad
 \tan\phi=\frac{2\beta\omega}{\omega^2-\omega_0^2}
\]
（$\phi$ 的象限由运动方程和电子电荷确定）。
平均从场吸收的功率等于平均阻尼耗散，
\[
 \overline P=\gamma\overline{\dot x^2}
 =\frac{\gamma a^2\omega^2}{2}
 =\boxed{\frac{\gamma e^2E_0^2\omega^2}
 {2[(\kappa-m_e\omega^2)^2+\gamma^2\omega^2]}}.
\]
对正频率求极值，在 $\omega=\omega_0$ 处
\[
 \boxed{\overline P_{\max}=\frac{e^2E_0^2}{2\gamma}
 =\frac{m_e}{4\beta}\left(\frac{eE_0}{m_e}\right)^2}.
\]
\excheck 此处是\emph{吸收功率}的极值频率；
阻尼振幅的最大点在一般情况下并非恰好 $\omega_0$。

\dispersionproblem{205}{221}
某些媒质可具有复折射率
$\tilde n=n+i\kappa$，也可具有纯虚折射率
$\tilde n=i\kappa$。写出平面波方程，
并解释复折射率的物理意义。
\exsolution
取时间依赖惯例
$E(x,t)=\Re\{E_0e^{i(\tilde kx-\omega t)}\}$，
$\tilde k=\tilde n\,\omega/c$。
若 $\tilde n=n+i\kappa$，
\[
 \boxed{E(x,t)=E_0 e^{-\kappa\omega x/c}
 \cos\left(\frac{n\omega x}{c}-\omega t\right)}.
\]
其中 $n$ 描述相位传播，
$\kappa>0$ 描述振幅的指数衰减，
光强沿 $x$ 按 $e^{-2\kappa\omega x/c}$ 衰减。
当 $\tilde n=i\kappa$ 时，
\[
 \boxed{E(x,t)=E_0e^{-\kappa\omega x/c}\cos\omega t}.
\]
介质内部不存在实波矢对应的传播相位，
仅有倏逝场；在无耗散、半无限介质的截止频段，
时间平均的法向净传播能流为零。
\excheck 复折射率的虚部正负要和所选时空相位惯例一致；
它不是第二个独立的相速度。

\dispersionproblem{206}{221}
用不同频率的无线电波探测稀薄等离子体，
发现真空波长大于
$\lambda_0=0.75\,\mathrm m$ 的波均被反射。
求自由电子数密度。
\exsolution
等离子体截止角频率为
$\omega_p=2\pi c/\lambda_0$；
利用 $\omega_p^2=n_0e^2/(\varepsilon_0m_e)$，
\[
 \boxed{n_0=
 \frac{4\pi^2\varepsilon_0m_ec^2}
 {e^2\lambda_0^2}
 \simeq2.0\times10^{15}\,\mathrm{m^{-3}}
 =2.0\times10^9\,\mathrm{cm^{-3}}}.
\]
\excheck 截止处 $\varepsilon_r(\omega_p)=0$；
题中是理想、无碰撞模型下的波不能穿透。

\dispersionproblem{207}{221}
从群速度定义 $u=d\omega/dk$ 推导瑞利公式，
并解释原图5.36中
在 $\lambda=\lambda'$ 处曲线 $v(\lambda)$ 的
切线与纵轴交点的物理意义。
\exsolution
因 $\omega=2\pi v/\lambda$，
$k=2\pi/\lambda$，
\[
 u=\frac{d(2\pi v/\lambda)/d\lambda}
 {d(2\pi/\lambda)/d\lambda}
 =\boxed{v-\lambda\frac{dv}{d\lambda}}.
\]
在 $\lambda'$ 处切线方程
$y=v(\lambda')+v'(\lambda')(\lambda-\lambda')$。
令 $\lambda=0$，截距为
$v(\lambda')-\lambda'v'(\lambda')$，
正是群速度 $u(\lambda')$。
\excheck 这只用到局部微分定义，
并不要求整个色散关系为线性。

\dispersionproblem{208}{221}
求以下各色散关系下群速度 $u$ 与相速度 $v$ 的比例：
(a) $v\propto1/\sqrt\lambda$；
(b) $v\propto k$；
(c) $v\propto1/\omega^2$。
\exsolution
(a) 写 $v=A\lambda^{-1/2}$，
由瑞利公式
$u=v-\lambda\,dv/d\lambda
=v+\tfrac12v$，故
$\boxed{u=\tfrac32v}$。
(b) 若 $v=Bk$，则 $\omega=vk=Bk^2$，
$u=d\omega/dk=2Bk$，
故 $\boxed{u=2v}$。
(c) 若 $v=C/\omega^2$，
又 $v=\omega/k$，
故 $k=\omega^3/C$，
$u=d\omega/dk=C/(3\omega^2)$，
从而 $\boxed{u=v/3}$。
\excheck 原书扫描第221页(a)项清晰为
$1/\sqrt{\lambda}$，PDF文字提取层漏掉根号；
若误读 $1/\lambda$ 会错误地得到 $u=2v$。

\dispersionproblem{209}{221--222}
某介质中电磁波满足
群速度与相速度乘积 $uv=c^2$，
求相对介电常数随频率的函数形式
$\varepsilon_r(\omega)$，假设磁导率约为1。
\exsolution
令相速度 $v=\omega/k$，
群速度 $u=d\omega/dk$，
则
\[
 \frac\omega k\frac{d\omega}{dk}=c^2,
 \quad \omega\,d\omega=c^2 k\,dk.
\]
积分得
$\omega^2=c^2k^2+C$。
由于
$n^2=\varepsilon_r=c^2k^2/\omega^2$，
\[
 \boxed{\varepsilon_r(\omega)=1+\frac{A}{\omega^2}},
 \qquad A=-C=\mathrm{const}.
\]
\excheck 若介质取无碰撞电子等离子体，
$A=-\omega_p^2$；仅从速度关系并不能判定常数符号。

\dispersionproblem{210}{222}
某介质在波长
$509,534,589\,\mathrm{nm}$ 的折射率依次为
$1.647,1.640,1.630$。
估算 $\lambda=534\,\mathrm{nm}$ 附近的
光相速度与群速度。
\exsolution
相速度为
\[
 \boxed{v=\frac c{1.640}\simeq1.83\times10^8\,\mathrm{m/s}}.
\]
在 $534\,\mathrm{nm}$ 附近用两侧点作差分，
\[
 \frac{dn}{d\lambda}\simeq
 \frac{1.630-1.647}{589-509}
 =-2.125\times10^{-4}\,\mathrm{nm^{-1}}.
\]
由 $v=c/n(\lambda)$ 和瑞利公式，
\[
 u=v-\lambda\frac{dv}{d\lambda}
 =\frac c n\left(1+\frac{\lambda}{n}\frac{dn}{d\lambda}\right),
\]
故
\[
 \boxed{u\simeq1.70\times10^8\,\mathrm{m/s}}.
\]
\excheck 正常色散下 $dn/d\lambda<0$，这里 $u<v$；
估计值随数值差分取法有小幅变化。

\dispersionproblem{211}{222}
一列任意形状的平面光波在某媒质中传播，
相速度与波长关系为
$v=a+b\lambda$，$a,b>0$。
证明任意波列的形状经过时间
$\tau=1/b$ 恢复（除整体平移）。
\exsolution
用波数 $k=2\pi/\lambda$，
\[
 \omega(k)=vk=ak+2\pi b.
\]
将任意波列作傅里叶叠加：
\[
 E(x,t)=\Re\!\int A(k)
 e^{i[kx-\omega(k)t]}\,dk
 =\Re\!\left\{
 e^{-i2\pi bt}
 \int A(k)e^{ik(x-at)}\,dk\right\}.
\]
所有谱分量共有的相位因子
$e^{-i2\pi bt}$ 在
$t=\tau=1/b$ 后恢复到1，
而积分只发生整体平移 $x\mapsto x-at$。
所以
\[
 \boxed{E(x,\tau)=E(x-a\tau,0)}.
\]
\excheck 群速度 $d\omega/dk=a$；
若只关心复包络形状，任意时刻都保持，
$t=1/b$ 是\emph{实场载波的共同相位}也复原的周期。

\dispersionproblem{212}{222}
自然光强度 $I_0$，依次通过交叉尼科耳棱镜，
中间夹一长 $l$ 的溶液管，沿管加纵向磁场 $H$。
溶液线吸收系数为 $\chi$，费尔德常数为 $V$。
求透射强度。
\exsolution
第一尼科耳棱镜后强度为 $I_0/2$。
沿管传播因吸收乘上 $e^{-\chi l}$，
并发生法拉第旋光 $\phi=VHl$。
再经与第一片交叉的尼科耳棱镜，
透射比例为 $\sin^2\phi$。故
\[
 \boxed{I=\frac{I_0}{2}e^{-\chi l}\sin^2(VHl)}.
\]
\excheck $H=0$ 时交叉检偏器应完全消光；
本题忽略界面反射，吸收系数作用于\emph{强度}。

\dispersionproblem{213}{222}
平面单色光强度 $I_0$ 垂直入射一块
两面强度反射率都为 $\rho$ 的平行平板。
计入多次内部反射，求
(a) 无吸收；
(b) 厚 $d$、线吸收系数 $\chi$ 时的总透射强度。
\exsolution
这里按书末采用各次透射光束\emph{强度非相干相加}
的近似（厚板或相位平均）。
设单程强度衰减因子 $s=e^{-\chi d}$。
第一条直通光线透射强度
$I_0(1-\rho)^2s$，
其后每多一次往返增加反射损耗 $\rho^2$
与两次穿行损耗 $s^2$。
故几何级数求和给
\[
 \boxed{I_{\mathrm{tr}}=
 I_0\frac{(1-\rho)^2s}{1-\rho^2s^2}}.
\]
(a) $s=1$ 时
\[
 \boxed{I_{\mathrm{tr}}=I_0\frac{1-\rho}{1+\rho}}.
\]
(b) $s=e^{-\chi d}$ 时
\[
 \boxed{I_{\mathrm{tr}}=
 I_0\frac{(1-\rho)^2e^{-\chi d}}
 {1-\rho^2e^{-2\chi d}}}.
\]
\excheck 对严格相干的单色平行光，
法布里--珀罗干涉使透射强度另依光程相位而变化；
本题无相位信息，书末所求为强度级数的相位平均模型。

\dispersionproblem{214}{222}
同一种均匀材料制成两片平板，厚度分别
$d_1=3.8\,\mathrm{mm}$ 与 $d_2=9.0\,\mathrm{mm}$。
对相同入射单色光的能流透过比例分别为
$\tau_1=0.84$ 和 $\tau_2=0.70$。
忽略二次反射，求材料的线吸收系数 $\chi$。
\exsolution
两片表面材料相同、入射角相同，
两片的表面共同透射因子 $C$ 可以约去：
\[
 \tau_1=C e^{-\chi d_1},\quad
 \tau_2=C e^{-\chi d_2}.
\]
由二式相除可得
\[
 \boxed{\chi=\frac{\ln(\tau_1/\tau_2)}{d_2-d_1}
 =\frac{\ln(0.84/0.70)}{0.52\,\mathrm{cm}}
 \simeq0.351\,\mathrm{cm^{-1}}}.
\]
\excheck 不能直接用 $-\ln\tau_i/d_i$，
那会把未知的两面反射损失误算为吸收。

\dispersionproblem{215}{222}
单色光依次穿过 $N=5$ 块相同的平行玻璃平板，
每片厚 $l=0.50\,\mathrm{cm}$。
每个表面的反射率 $\rho=0.050$，
五片总透射能流比 $\tau=0.55$。
忽略二次反射，求玻璃的线吸收系数 $\chi$。
\exsolution
每片有两面，反射导致的总透射因子是
$(1-\rho)^{2N}$；
总材料厚度为 $Nl$，吸收因子为 $e^{-\chi Nl}$。
所以
\[
 \tau=(1-\rho)^{2N}e^{-\chi Nl},
 \quad
 \boxed{\chi=\frac1{Nl}
 \ln\frac{(1-\rho)^{2N}}{\tau}}.
\]
代入给定量：
\[
 \boxed{\chi=\frac1{2.50\,\mathrm{cm}}
 \ln\frac{0.95^{10}}{0.55}
 \simeq0.034\,\mathrm{cm^{-1}}}.
\]
\excheck 材料内吸收与界面反射分别计数；
计算所得 $\chi>0$ 与题设相符。
